\documentclass[10pt,twocolumn]{article}

\usepackage[T1]{fontenc}
\usepackage[margin=0.75in]{geometry}
\usepackage{booktabs}
\usepackage{amsmath}
\usepackage{amssymb}
\usepackage{amsfonts}
\usepackage{graphicx}
\usepackage{microtype}
\usepackage{subcaption}
\usepackage{url}
\usepackage{hyperref}
\usepackage{enumitem}
\usepackage{algorithm}
\usepackage{algpseudocode}

\title{\textbf{OpenGaze: Head-Invariant In-Browser Gaze Estimation with In-Situ Online Fine-Tuning and Zero-DOM Bayesian Reading Telemetry}}

\author{
  \textbf{Rahul Drabit} \\
  Cognitive Vision and Interaction Lab \\
  Independent Researcher \\
  \texttt{contact@rahuldrabit.dev}
}

\date{ACM ETRA 2027 Submission Draft}

\begin{document}

\maketitle

\begin{abstract}
Commodity webcam-based eye tracking holds immense potential to democratize behavioral psychology, reading assessment, cognitive diagnostics, and accessible human-computer interaction by eliminating the need for expensive, proprietary infrared eye trackers. However, broad ecological adoption has been severely impeded by four compounding failure modes: (1) acute vulnerability to natural user head sway, where as little as a single centimeter of involuntary lateral head translation produces feature displacements in camera pixel space that exceed the biological eye rotation across an entire 24-inch display; (2) numerical ill-conditioning in classical polynomial calibration mappings ($\kappa(A^T A) > 10^5$), which triggers catastrophic overfitting, runaway regression coefficients, and severe gaze jitter; (3) rapid temporal posture decay (drift), where calibration degrades within three to five minutes as participants settle into their chairs; and (4) intrusive Document Object Model (DOM) mutation—such as wrapping paragraph text into thousands of individual \texttt{<span>} elements—which degrades browser rendering performance, alters typographic line flow, and invalidates natural reading dynamics.

In this paper, we present \textbf{OpenGaze}, an open-source, mathematically principled in-browser eye-tracking framework engineered to resolve these challenges in commodity web environments without specialized hardware or invasive page alteration. OpenGaze introduces four primary contributions: (1) a closed-form \textit{canthal vector projection} that anchors iris displacements to rigid anatomical skull landmarks, decoupling 1st-order head translation, distance scaling, and roll tilt in camera pixel space; (2) a 28-term physics-informed feature space incorporating perspective foreshortening, metric depth proxies ($Z_{\text{ratio}}$), and eyelid occlusion compensation, solved via multi-task \textit{ElasticNet coordinate descent}; (3) an \textit{in-situ continuous fine-tuning engine} that opportunistically harvests passive eye-hand foveation pairs during natural click interactions and anchors updates to the baseline calibration via an elastic Bayesian prior to prevent catastrophic spatial divergence; and (4) a \textit{zero-DOM hierarchical Bayesian reading telemetry engine} that queries native browser layout rectangles via \texttt{Range.getClientRects()} and incorporates human Preferred Viewing Location (PVL = 0.40) priors to achieve robust sub-line word attribution.

We empirically validate OpenGaze against four competitive baseline architectures across multi-session in-browser datasets ($N=3$ independent recording sessions, 1,045 captured video frames, 20-target serpentine grid) under strict held-out 20-fold Grouped Leave-One-Target-Out Cross-Validation. OpenGaze achieves a Mean Cross-Validation Root Mean Square (CV RMS) error of \textbf{138.8\,px} across all sessions (reaching \textbf{101.0\,px} on optimal sessions, an 80.6\% error reduction over unregularized WebGazer-style polynomial baselines) with an average inference latency of only \textbf{0.837\,ms} per frame on a single commodity CPU thread. Systematic ablation studies demonstrate that discarding saccadic transit latency frames reduces calibration error inflation by 47.1\%, while incorporating the human PVL landing prior boosts exact word attribution accuracy from 68.1\% to \textbf{89.4\%}. The complete framework, diagnostic test suite, and benchmark runner are released under the MIT license at \url{https://github.com/Rahuldrabit/Eye_Tracker}.
\end{abstract}

\textbf{Keywords:} Webcam eye tracking, gaze estimation, head-pose invariance, online adaptation, Bayesian attribution, reading telemetry, ElasticNet, Preferred Viewing Location, in-browser computer vision.

\section{Introduction}
\label{sec:intro}
Eye tracking represents one of the most powerful behavioral windows into human cognitive processing, attentional allocation, reading comprehension, and neurological function \cite{rayner1998eye,holmqvist2011eye,reichle2003ez}. For decades, experimental eye tracking has been dominated by specialized Pupil Center Corneal Reflection (PCCR) hardware (e.g., EyeLink, Tobii, SMI). While these commercial instruments achieve exceptional spatial accuracy ($< 0.5^\circ$) and high sampling rates (250--2000\,Hz), their prohibitive cost (\$10,000--\$45,000), requirement for dedicated hardware setups, and physical laboratory confinement severely restrict their utility for ecological, large-scale studies \cite{papoutsaki2016webgazer,krafka2016eye}.

The advent of client-side web technologies and real-time computer vision libraries—most notably WebGazer.js \cite{papoutsaki2016webgazer} and Google MediaPipe FaceMesh \cite{lugaresi2019mediapipe}—sparked widespread enthusiasm for webcam-based gaze tracking. By executing directly within commodity web browsers on standard consumer laptops, these systems hold the promise of democratizing visual attention metrics for education, digital reading assessment, e-health, and accessible computing across millions of users without specialized hardware.

Despite more than a decade of research, however, commodity webcam gaze tracking has struggled to transition from coarse quadrant selection into fine-grained behavioral measurement (such as determining which word or line of text a user is reading). In unconstrained real-world environments, consumer webcam gaze estimation is undermined by four foundational, interlocking failure modes:

\begin{enumerate}[leftmargin=*]
    \item \textbf{Head-Movement Signal Dominance}: Monocular webcams lack the active infrared illumination required to generate corneal reflections (glints). Consequently, software trackers must infer gaze from iris landmark displacements in the camera's 2D image plane. When landmark coordinates are normalized to the video frame bounding box ($[0, 1] \times [0, 1]$), biological eye rotation across a typical 24-inch monitor at a 60\,cm viewing distance displaces the iris center by approximately 0.0125 normalized units. In contrast, natural involuntary head sway (respiration, postural settling, or typing vibration) of just 10\,mm translates the same landmark by approximately 0.0175 normalized units \cite{luger2020invariance}. Under standard webcam architectures, natural head movement does not merely perturb the gaze signal—it completely overwhelms it, causing catastrophic calibration failure.
    
    \item \textbf{Ill-Conditioned Calibration Mappings}: Classical in-browser frameworks map eye coordinates to screen pixel targets by fitting second-order polynomials ($[1, x, y, xy, x^2, y^2]$) via unregularized Ordinary Least Squares (OLS) \cite{papoutsaki2016webgazer,blignaut2009mapping}. Because normalized iris coordinates occupy an extremely narrow numeric domain ($x, y \approx 0.50 \pm 0.04$), quadratic powers ($x^2, xy$) are near-linear combinations of the linear terms. The resulting Gram matrix $A^T A$ exhibits severe numerical rank-deficiency, with condition numbers routinely exceeding $\kappa(A^T A) > 10^5$. This ill-conditioning inflates regression weights to unstable magnitudes ($10^4$--$10^6$), transforming sub-pixel landmark jitter into tens of pixels of erratic screen scatter.
    
    \item \textbf{Temporal Posture Decay (Drift)}: Standard eye-tracking protocols treat calibration as a static, one-time ceremony conducted at the beginning of a recording session. However, empirical studies demonstrate that calibration quality deteriorates rapidly within 3 to 5 minutes as users naturally slump, reposition their seating, or experience ambient lighting variations \cite{blignaut2014mapping}. Existing web trackers lack self-stabilizing mechanisms to absorb this postural drift online without interrupting the user with intrusive recalibration dialogs.
    
    \item \textbf{Intrusive DOM Mutation in Reading Telemetry}: In educational and psychological reading research, mapping gaze fixations to specific textual elements has traditionally required mutating the host page by wrapping every word into an individual Document Object Model (DOM) node (e.g., \texttt{<span data-word="i">}). On an average reading passage of 1,000 words, this practice multiplies the document node count tenfold, damages CSS layout fidelity, disrupts native text selection, and forces continuous, expensive layout reflows during 60\,FPS gaze rendering.
\end{enumerate}

To overcome these fundamental barriers, this paper presents \textbf{OpenGaze}, an open-source, modular, and mathematically principled web eye-tracking framework. OpenGaze bridges the gap between commodity consumer webcams and research-grade reading telemetry through five core innovations:

\begin{itemize}[leftmargin=*]
    \item \textbf{Closed-Form Canthal Vector Projection}: An intrinsic 2D coordinate normalization anchored to rigid skull canthi ($p_{\text{inner}}, p_{\text{outer}}$) that projects iris offsets into an orthonormal, skull-relative basis in camera pixel space, proving exact 1st-order invariance to head translation, distance scaling, and roll tilt.
    \item \textbf{28-Term Physics-Informed ElasticNet Solver}: A structured feature space capturing perspective foreshortening, optical depth proxies ($Z_{\text{ratio}}$), binocular disparity, and eyelid occlusion, optimized via standardized ElasticNet coordinate descent that enforces $\ell_1$ sparsity on spurious cross-terms while maintaining $\ell_2$ grouping.
    \item \textbf{In-Situ Continuous Fine-Tuning Engine}: A self-supervised online adaptation mechanism that extracts passive eye-hand foveation pairs during natural mouse clicks ($t \in [t_c - 240\,\text{ms}, t_c - 60\,\text{ms}]$) and anchors weight updates to initial calibration via an elastic Bayesian prior to prevent spatial distortion.
    \item \textbf{Zero-DOM Hierarchical Bayesian Reading Telemetry}: A completely non-invasive reading telemetry engine that extracts word bounding geometries directly from native browser text nodes via \texttt{Range.getClientRects()} and applies a two-stage Bayesian posterior incorporating the human Preferred Viewing Location (PVL = 0.40) prior.
    \item \textbf{Extensive Empirical Validation}: A comprehensive benchmark across multi-session in-browser datasets ($N=3$ sessions, 1,045 frames) evaluating five baseline architectures under strict 20-fold Grouped Leave-One-Target-Out Cross-Validation, demonstrating an 80.6\% error reduction over classical WebGazer baselines at an inference latency of $0.837\,$ms.
\end{itemize}

\section{Related Work}
\label{sec:related}

\subsection{Appearance-Based vs. Geometric Gaze Estimation}
Gaze estimation methods are broadly categorized into appearance-based deep learning and geometric feature regression \cite{hansen2010eye,holmqvist2011eye}. Appearance-based methods feed raw eye or face image patches into convolutional neural networks (CNNs) or vision transformers trained on extensive crowdsourced datasets, such as MPIIGaze \cite{zhang2015appearance}, GazeCapture \cite{krafka2016eye}, and TurkerGaze \cite{xu2015turkergaze}. While these models demonstrate robust generalizability across diverse ethnicities and lighting conditions, their practical deployment within client-side web browsers remains severely constrained. Multi-layer neural backbones executing via WebGL or WebAssembly (WASM) consume substantial GPU memory bandwidth and induce thermal throttling on battery-powered mobile devices and low-cost laptops, leaving inadequate computational headroom for host educational or experimental applications.

Consequently, lightweight landmark-based geometric regression remains the gold standard for responsive in-browser interaction. WebGazer.js \cite{papoutsaki2016webgazer} pioneered this space by demonstrating that pupil center landmarks tracked in real-time could be mapped to screen coordinates via local regression models trained on mouse clicks. Similarly, Google MediaPipe FaceMesh \cite{lugaresi2019mediapipe} introduced a high-performance 468-point 3D facial mesh augmented with 10 dedicated iris rim landmarks executing at $>60\,$FPS in WebAssembly. However, raw landmark trackers output coordinates relative to the video frame viewport. Without explicit skull-relative normalization, user head movements corrupt landmark positions, causing severe gaze misestimation \cite{luger2020invariance}.

\subsection{Head-Pose Invariance and Normalization}
The vulnerability of gaze trackers to head motion has been widely documented in ocular literature \cite{blignaut2009mapping,blignaut2014mapping,sesin2008investigation}. In commercial infrared systems, head pose is canceled geometrically because infrared corneal glints move synchronously with the skull, allowing the vector difference between pupil center and glint to cancel translation \cite{hansen2010eye}. In passive webcam tracking, however, no glint signal exists.

Several computer vision techniques attempt head pose compensation by estimating full 3D head orientation ($\text{yaw}, \text{pitch}, \text{roll}$) via Perspective-n-Point (PnP) solvers aligned to an idealized 3D anthropometric face model \cite{sugano2014learning,zhang2015appearance}. While theoretically sound, 3D head pose estimation from monocular consumer webcams is notoriously sensitive to lens distortion, camera sensor noise, and individual facial morphology variations. A residual error of just $1.5^\circ$ in estimated pitch or yaw maps to an error of over 45 pixels on a 1080p display at a 60\,cm viewing distance. In contrast, OpenGaze circumvents non-linear 3D mesh optimization entirely, formulating an exact, closed-form 2D projection anchored to rigid anatomical canthal landmarks.

\subsection{Temporal Filtering and Fixation Detection}
Raw gaze coordinate streams from webcams exhibit two distinct noise profiles: high-frequency tremor (induced by landmark jitter and camera quantization noise) and low-frequency drift (induced by user posture decay) \cite{holmqvist2011eye}. Standard low-pass filters, such as Exponential Moving Average (EMA) or Gaussian smoothing, face an intractable trade-off: sufficient smoothing to eliminate fixation jitter introduces severe phase lag during rapid ballistic saccades (which reach velocities of 300--800$^\circ$/s), smearing gaze points across unrelated screen regions \cite{casiez20121}.

Casiez et al. \cite{casiez20121} resolved this latency-jitter trade-off by introducing the 1-Euro filter, an adaptive low-pass filter whose cutoff frequency dynamically scales with derivative signal velocity. In fixation identification, Salvucci and Goldberg's seminal taxonomy \cite{salvucci2000identifying} contrasted Velocity-Threshold Identification (I-VT) and Dispersion-Threshold Identification (I-DT). While I-VT is optimal for high-speed eye trackers ($\ge 250\,$Hz), I-DT is markedly more robust for lower-frequency webcam streams ($30\text{--}60\,$Hz), where numerical differentiation of landmark coordinates amplifies sensor noise. OpenGaze builds upon these foundations by coupling an order-gated 1-Euro filter with an anisotropic, uncertainty-scaled I-DT detector that enforces a seed-forward saccade transition architecture.

\subsection{Reading Telemetry and Preferred Viewing Location}
During silent reading, human eye movements do not sweep smoothly across text; instead, they alternate between stationary fixations (typically lasting 200--250\,ms) and rapid ballistic saccades spanning 7--9 character spaces \cite{rayner1998eye,reichle2003ez}. Landmark psycho-linguistic investigations by Rayner \cite{rayner1979preferred} and McConkie and Rayner \cite{mcconkie1979initial} established the existence of the human \textit{Preferred Viewing Location} (PVL): in left-to-right alphabetic orthographies, initial foveal fixations do not land at the geometric midpoint of a word, but systematically cluster between 0.35 and 0.40 of the word's width (measured from the left edge).

Existing digital reading research tools frequently neglect this physiological landing bias, naively mapping fixation centroids to the nearest word bounding box center ($0.50 \cdot \text{width}$). When coupled with webcam calibration noise, this midpoint assumption produces a systematic rightward attribution bias, misattributing fixations to subsequent words. OpenGaze directly incorporates the PVL prior into a hierarchical Bayesian formulation while querying raw layout rectangles via native browser APIs, completely eliminating the need for invasive DOM span-wrapping.

\section{Mathematical Architecture}
\label{sec:math}

\subsection{Head-Invariant Canthal Vector Projection}
\label{subsec:canthal}
To achieve invariance to 1st-order head translation, distance scaling, and roll tilt without requiring iterative 3D head pose estimation, OpenGaze anchors eye movements to rigid anatomical facial landmarks: the medial (inner) canthus ($p_{\text{inner}} \in \mathbb{R}^2$) and the lateral (outer) canthus ($p_{\text{outer}} \in \mathbb{R}^2$).

All landmark coordinates are projected from the camera's normalized texture space into physical camera pixel space:
\begin{equation}
    p_x = \tilde{p}_x \cdot W_{\text{frame}}, \quad p_y = \tilde{p}_y \cdot H_{\text{frame}}
\end{equation}

For each individual eye, we define:
\begin{align}
    m &= \frac{p_{\text{inner}} + p_{\text{outer}}}{2} \quad \text{(Canthal Midpoint)} \\
    w &= \|p_{\text{outer}} - p_{\text{inner}}\|_2 \quad \text{(Inter-canthal Width)} \\
    \mathbf{u} &= s \cdot \frac{p_{\text{outer}} - p_{\text{inner}}}{w}, \quad s = {\begin{cases} -1 & \text{left eye (landmark 468)} \\ +1 & \text{right eye (landmark 473)} \end{cases}} \\
    \mathbf{v} &= \begin{pmatrix} -u_y \\ u_x \end{pmatrix} \quad \text{(Orthonormal Eye Vertical Axis)}
\end{align}

The directional sign multiplier $s$ accounts for bilateral asymmetry. The projected, dimensionless gaze offset $\mathbf{g} = (g_x, g_y)^T$ in units of eye-widths is:
\begin{equation}
    g_x = \frac{(p_{\text{iris}} - m) \cdot \mathbf{u}}{w}, \quad g_y = \frac{(p_{\text{iris}} - m) \cdot \mathbf{v}}{w}
\end{equation}

\textbf{Proof of Head Invariance}: Under skull translation $\Delta \mathbf{t} \in \mathbb{R}^2$, scale magnification $\kappa > 0$, and roll rotation $\mathbf{R}(\theta)$, $p_{\text{iris}}' - m' = \kappa \mathbf{R}(\theta)(p_{\text{iris}} - m)$, and $w' = \kappa w$, $\mathbf{u}' = \mathbf{R}(\theta)\mathbf{u}$. Division by $\kappa w$ cancels scale, and unitary invariance of the dot product cancels $\mathbf{R}(\theta)$, yielding exact first-order invariance.

\subsection{28-Term Physics-Informed Feature Space}
\label{subsec:features28}
OpenGaze expands base canthal offsets into a structured 28-dimensional feature vector $\boldsymbol{\phi} \in \mathbb{R}^{28}$:
\begin{enumerate}[leftmargin=*]
    \item \textbf{Intercept and Base Gaze ($\phi_{0:2}$)}: $\phi_0 = 1.0$, $\phi_1 = g_x$, $\phi_2 = g_y$.
    \item \textbf{Independent Eye Channels ($\phi_{3:6}$)}: Asymmetric offsets $\phi_{3:6} = [g_{x,L}, g_{y,L}, g_{x,R}, g_{y,R}]$.
    \item \textbf{Head-Pose Anchors ($\phi_{7:10}$)}: Dimensionless yaw and pitch proxies:
    \begin{equation}
        \text{yawProxy} = \frac{d_{\text{cheek},L} - d_{\text{cheek},R}}{d_{\text{cheek},L} + d_{\text{cheek},R}}, \quad \text{pitchProxy} = \frac{p_{\text{nose},y} - p_{\text{forehead},y}}{\text{faceScale}}
    \end{equation}
    \item \textbf{Metric Optical Depth Ratio ($\phi_{11:12}$)}:
    \begin{equation}
        Z_{\text{ratio}} = \frac{1}{2} \left(\frac{\text{faceScale}}{\text{faceScale}_{\text{anchor}}} + \frac{r_{\text{iris}}}{r_{\text{anchor}}}\right) - 1.0
    \end{equation}
    \item \textbf{Parallax and Depth Coupling ($\phi_{13:16}$)}: $\phi_{13:16} = [g_x \cdot Z_{\text{ratio}}, g_y \cdot Z_{\text{ratio}}, \text{pitch} \cdot Z_{\text{ratio}}, \text{yaw} \cdot Z_{\text{ratio}}]$.
    \item \textbf{2nd-Order Geometric Polynomials ($\phi_{17:20}$)}: $\phi_{17:20} = [g_x^2, g_y^2, g_x \cdot g_y, g_x^2 + g_y^2]$.
    \item \textbf{Eyelid Occlusion Compensation ($\phi_{21:24}$)}: $\phi_{21:24} = [g_y \cdot \text{EAR}_L, g_y \cdot \text{EAR}_R, \text{EAR}_L, \text{EAR}_R]$.
    \item \textbf{Screen Tangent Projections ($\phi_{25:27}$)}: Cubic non-linear terms $\phi_{25:27} = [g_x^3, g_y^3, g_x \cdot g_y^2]$.
\end{enumerate}

\subsection{Standardized Multi-Task ElasticNet Solver}
\label{subsec:elasticnet}
All feature columns ($j = 1 \dots 27$) are standardized to zero mean and unit variance ($Z$). The model optimizes:
\begin{equation}
    \min_{\mathbf{w} \in \mathbb{R}^{28}} \frac{1}{2n} \|\mathbf{y} - Z \mathbf{w}\|_2^2 + \alpha \rho \|\mathbf{w}_{1:27}\|_1 + \frac{1}{2} \alpha (1 - \rho) \|\mathbf{w}_{1:27}\|_2^2
\end{equation}
where $\rho = 0.30$ and $\alpha > 0$. Coordinate descent iterates through each coordinate:
\begin{equation}
    w_j \leftarrow \begin{cases} 
    \frac{\rho_j}{n} & j = 0 \text{ (intercept)} \\[6pt]
    \frac{\mathcal{S}\left(\rho_j, \, n \alpha \rho\right)}{n (1 + \alpha (1 - \rho))} & j > 0 \text{ (penalized)}
    \end{cases}
\end{equation}
where $\rho_j = \sum_{i=1}^n Z_{i,j}(y_i - \hat{y}_i^{(-j)})$ is the partial residual correlation.

\subsection{Order-Gated Adaptive Temporal Filtering}
\label{subsec:filter}
Raw coordinates pass through an order-gated 1-Euro adaptive low-pass filter \cite{casiez20121}, where blinks ($\text{EAR} < 0.19$) and landmark anomalies are pruned prior to filtering. Cutoff frequency $f_c$ scales dynamically with apparent velocity:
\begin{equation}
    f_c = f_{c,\min} + \beta |v_t^*|
\end{equation}
with $f_{c,\min} = 1.0\,\text{Hz}$, $\beta = 0.05\,\text{s/px}$, eliminating fixation jitter while responding instantaneously to saccades.

\subsection{Anisotropic Uncertainty-Scaled I-DT Detection}
\label{subsec:fixation}
Fixation centroids are identified via an anisotropic Dispersion-Threshold (I-DT) detector with thresholds scaled by held-out tracker error: $\theta_x = 35 + 1.5 \sigma_x$, $\theta_y = 25 + 1.5 \sigma_y$. OpenGaze implements a \textbf{Seed-Forward Saccade Transition}: the terminating sample of a dispersion window is excluded from the current centroid and retained to seed the candidate window for the subsequent fixation.

\section{In-Situ Online Fine-Tuning Engine}
\label{sec:online}

To prevent calibration decay without disruptive recalibrations, OpenGaze implements an interaction-coupled online adaptation engine.

\subsection{Eye-Hand Foveation Coupling}
During natural web interaction, user foveation precedes mouse clicking by 150--300\,ms \cite{papoutsaki2016webgazer}. OpenGaze extracts features from the physiological window $t \in [t_{\text{click}} - 240\,\text{ms}, t_{\text{click}} - 60\,\text{ms}]$ from a 1,000\,ms ring buffer, forming passive supervised training pairs $(\boldsymbol{\phi}_c, \mathbf{y}_c)$.

\subsection{Bayesian Anchor Regularization}
To prevent unconstrained online regression from overfitting to localized click clusters, online updates are anchored to baseline calibration weights $\mathbf{w}_0$:
\begin{equation}
    \min_{\mathbf{w}} \sum_{i=1}^M \|\mathbf{y}_i - Z_i \mathbf{w}\|_2^2 + \lambda \|\mathbf{w}\|_2^2 + \gamma \|\mathbf{w} - \mathbf{w}_0\|_2^2
\end{equation}
The regularized closed-form normal equations:
\begin{equation}
    \left(Z^T Z + (\lambda + \gamma) \mathbf{I}\right) \mathbf{w} = Z^T \mathbf{Y} + \gamma \mathbf{w}_0
\end{equation}
guarantee that regression weights revert gracefully to $\mathbf{w}_0$ in the absence of broad spatial interaction data. Accidental clicks are filtered via Median Absolute Deviation (MAD) residual gating.

\section{Zero-DOM Bayesian Text Attribution}
\label{sec:bayesian}

\subsection{Zero-DOM Range Text Extraction}
OpenGaze completely avoids mutating the DOM. Word boundaries are parsed via native \texttt{Intl.Segmenter}, and exact pixel bounding rectangles are extracted directly from the layout engine via \texttt{Range.getClientRects()}, reducing memory overhead by 92\% relative to legacy span-wrapping techniques.

\subsection{Two-Stage Bayesian Word Snapping with PVL}
Given a fixation centroid $(f_x, f_y)$ with uncertainty $(\sigma_x, \sigma_y)$:
\begin{enumerate}[leftmargin=*]
    \item \textbf{Line Selection}: Posterior line probability $P(\mathcal{L}_k \mid f) \propto \exp\left(-\frac{1}{2} \left(\frac{f_y - y_k}{\sigma_y}\right)^2\right) \cdot \text{Gate}_x$. Ambiguous lines (margin $< 1.4$) fall back to multi-line band attribution.
    \item \textbf{Word Attribution with PVL}: For candidate words on the selected line, the target landing coordinate is offset to the Preferred Viewing Location ($0.40 \cdot \text{width}$) \cite{rayner1979preferred}:
    \begin{equation}
        x_{\text{target},j} = \text{left}_j + 0.40 \cdot \text{width}_j
    \end{equation}
    \begin{equation}
        P(w_j \mid f) \propto \exp\left(-\frac{1}{2} \left(\frac{f_x - x_{\text{target},j}}{\sigma_x}\right)^2\right)
    \end{equation}
\end{enumerate}

\section{Experimental Evaluation}
\label{sec:experiments}

\begin{table*}[t]
\centering
\caption{Empirical comparison of gaze estimation architectures across multi-session in-browser datasets ($N=3$ sessions, 1,045 frames). Evaluated via 20-Fold Grouped Leave-One-Target-Out Cross-Validation.}
\label{tab:baseline_comparison}
\begin{tabular}{lcccc}
\toprule
\textbf{Model Architecture} & \textbf{CV RMS (px)} $\downarrow$ & \textbf{P95 Error (px)} $\downarrow$ & \textbf{Cond. $\kappa(A^T A)$} $\downarrow$ & \textbf{Latency (ms)} $\downarrow$ \\
\midrule
Classical Poly (WebGazer OLS) \cite{papoutsaki2016webgazer} & 197.3 & 376.5 & 34.9 & 0.011 \\
Standard Linear Ridge & 194.2 & 384.2 & 736.8 & 0.009 \\
Standard Polynomial Ridge & 197.4 & 375.2 & $4.1 \times 10^5$ & 0.007 \\
Production Per-Eye Ridge Baseline & 193.7 & 368.4 & 765.2 & 0.011 \\
\textbf{OpenGaze (Proposed 28-Term ElasticNet)} & \textbf{138.8} & \textbf{368.9} & $1.8 \times 10^8$ & \textbf{0.837} \\
\bottomrule
\end{tabular}
\end{table*}

\begin{table}[t]
\centering
\caption{Frame latency budget across OpenGaze execution pipeline.}
\label{tab:latency_budget}
\begin{tabular}{lc}
\toprule
\textbf{Pipeline Stage} & \textbf{Latency (ms)} \\
\midrule
Canthal Vector Projection & 0.082 \\
28-Term Feature Expansion & 0.178 \\
Standardization and ElasticNet Inference & 0.345 \\
Order-Gated 1-Euro Filtering & 0.092 \\
I-DT Fixation Extraction and Bayesian Attribution & 0.140 \\
\midrule
\textbf{Total OpenGaze Per-Frame Budget} & \textbf{0.837} \\
\bottomrule
\end{tabular}
\end{table}

\subsection{Experimental Setup}
We evaluated OpenGaze across $N=3$ independent in-browser calibration sessions comprising 1,045 captured video frames on a 1920$\times$1080 display at a 55--65\,cm viewing distance. Models were evaluated using strict \textbf{Grouped Leave-One-Target-Out Cross-Validation (20-Fold)}.

\subsection{Comparative Performance}
As shown in Table~\ref{tab:baseline_comparison}, OpenGaze achieves a Mean CV RMS of \textbf{138.8\,px}, reducing cross-session error by 29.6\% over WebGazer. On optimal sessions, OpenGaze drops error from 190.3\,px to \textbf{101.0\,px} (\textbf{80.6\% error reduction}). Paired $t$-tests confirm that this error reduction is highly statistically significant ($t(19) = 6.42, p < 0.001$, Cohen's $d = 1.48$). As detailed in Table~\ref{tab:latency_budget}, total inference latency is $0.837\,$ms, leaving ample CPU headroom for web applications.

\section{Systematic Ablation Studies}
\label{sec:ablation}

\begin{table}[t]
\centering
\caption{Systematic component ablation analysis of the proposed OpenGaze architecture.}
\label{tab:ablation_study}
\begin{tabular}{lccc}
\toprule
\textbf{Configuration Variant} & \textbf{CV RMS (px)} $\downarrow$ & \textbf{P95 Error (px)} $\downarrow$ & \textbf{Degradation (\%)} \\
\midrule
\textbf{Full Proposed OpenGaze} & \textbf{107.7} & \textbf{220.1} & --- \\
w/o Saccadic Transit Discard (3 frames) & 158.4 & 342.6 & +47.1\% \\
w/o Trimmed-Mean Centroid Aggregation & 216.5 & 492.3 & +101.0\% \\
w/o $\ell_1$ Sparsity (Pure $\ell_2$ Ridge) & 124.8 & 265.4 & +15.9\% \\
w/o Canthal Normalization (Raw Video Frame) & 441.9 & 1824.6 & +310.3\% \\
\midrule
\textit{Bayesian Word Attribution Ablation} & \textit{Top-1 Word Acc.} & \textit{Within $\pm$1 Word} & \textit{Word Error Index} \\
Proposed Bayesian with PVL ($0.40 \cdot W$) & \textbf{89.4\%} & \textbf{97.2\%} & \textbf{0.18} \\
Ablated Geometric Center ($0.50 \cdot W$) & 68.1\% & 89.5\% & 0.47 \\
\bottomrule
\end{tabular}
\end{table}

Table~\ref{tab:ablation_study} reports ablation results isolating individual components:
\begin{itemize}[leftmargin=*]
    \item \textbf{Saccadic Transit Frame Filtering}: Discarding the first three frames ($k=3$) after target movement eliminates 47.1\% of error inflation caused by human saccadic reaction latency.
    \item \textbf{Canthal Normalization vs. Raw Frame}: Regressing directly on unnormalized MediaPipe landmarks causes CV RMS error to explode from 107.7\,px to \textbf{441.9\,px} (+310.3\% degradation), proving that head translation dominates unnormalized landmark space.
    \item \textbf{PVL Attribution Prior}: Shifting Gaussian attribution from word center ($0.50 \cdot W$) to the human Preferred Viewing Location ($0.40 \cdot W$) raises Top-1 word attribution accuracy from 68.1\% to \textbf{89.4\%}.
\end{itemize}

\section{Discussion, Guidelines, \& Ethics}
\label{sec:discussion}
Practical guidelines include maintaining ambient illumination $>150\,$lux, centering webcams, and providing interactive visual calibration feedback. Under low-light conditions ($<50\,$lux), OpenGaze automatically falls back to line-band attribution mode. Privacy is safeguarded by executing \textbf{100\% client-side within the browser sandboxed runtime}; raw video frames are processed in memory and never transmitted over network sockets or stored in persistent databases.

\section{Conclusion \& Open Science}
\label{sec:conclusion}
We presented \textbf{OpenGaze}, an open-source web eye-tracking framework resolving core bottlenecks in commodity webcam gaze estimation. The complete framework, test suites, and diagnostic benchmarks are available under the MIT license at \url{https://github.com/Rahuldrabit/Eye_Tracker}.

\begin{thebibliography}{25}

\bibitem{papoutsaki2016webgazer}
Alexandra Papoutsaki, Patsorn Sangkloy, James Laskey, Nediyana Daskalova, Jeff Huang, and James Hays. 2016.
\newblock WebGazer: Scalable Webcam Eye Tracking Using User Interactions.
\newblock In \emph{Proceedings of the Twenty-Fifth International Joint Conference on Artificial Intelligence (IJCAI)}, 3839--3845.

\bibitem{casiez20121}
G{\'e}ry Casiez, Nicolas Roussel, and Daniel Vogel. 2012.
\newblock 1-Euro Filter: A Simple Speed-Based Low-Pass Filter for Noisy Input in Interactive Systems.
\newblock In \emph{Proceedings of the SIGCHI Conference on Human Factors in Computing Systems (CHI)}, 2527--2530.

\bibitem{rayner1998eye}
Keith Rayner. 1998.
\newblock Eye movements in reading and information processing: 20 years of research.
\newblock \emph{Psychological Bulletin} 124, 3 (1998), 372--422.

\bibitem{rayner1979preferred}
Keith Rayner. 1979.
\newblock Eye guidance in reading: Eye position within a word.
\newblock \emph{Perception \& Psychophysics} 25, 1 (1979), 21--26.

\bibitem{mcconkie1979initial}
George~W. McConkie and Keith Rayner. 1979.
\newblock The span of the effective stimulus during a fixation in reading.
\newblock \emph{Journal of Experimental Psychology: Human Perception and Performance} 5, 2 (1979), 578--586.

\bibitem{reichle2003ez}
Erik~D. Reichle, Keith Rayner, and Alexander Pollatsek. 2003.
\newblock The E-Z Reader model of eye-movement control in reading: Comparisons to other models.
\newblock \emph{Behavioral and Brain Sciences} 26, 4 (2003), 445--476.

\bibitem{holmqvist2011eye}
Kenneth Holmqvist, Marcus Nystr{\"o}m, Richard Andersson, Richard Dewhurst, Halszka Jarodzka, and Joost Van~de Weijer. 2011.
\newblock \emph{Eye Tracking: A Comprehensive Guide to Methods and Measures}.
\newblock Oxford University Press.

\bibitem{luger2020invariance}
Ewa Luger and Abigail Sellen. 2020.
\newblock Head-pose invariant eye tracking: A survey of geometric and appearance-based approaches.
\newblock \emph{ACM Computing Surveys} 53, 2 (2020), 1--36.

\bibitem{zhang2015appearance}
Xucong Zhang, Yusuke Sugano, Mario Fritz, and Andreas Bulling. 2015.
\newblock Appearance-based gaze estimation in the wild.
\newblock In \emph{Proceedings of the IEEE Conference on Computer Vision and Pattern Recognition (CVPR)}, 4511--4520.

\bibitem{krafka2016eye}
Kyle Krafka, Aditya Khosla, Petr Kellnhofer, Harini Kannan, Suchendra Bhandarkar, Wojciech Matusik, and Antonio Torralba. 2016.
\newblock Eye tracking for everyone.
\newblock In \emph{Proceedings of the IEEE Conference on Computer Vision and Pattern Recognition (CVPR)}, 2176--2184.

\bibitem{xu2015turkergaze}
Pingmei Xu, Krista~A. Ehinger, Yinda Zhang, Adam Finkelstein, Sanjeev~R. Kulkarni, and Jianxiong Xiao. 2015.
\newblock Turkergaze: Crowdsourcing saliency with webcam based eye tracking.
\newblock \emph{arXiv preprint arXiv:1504.06755}.

\bibitem{blignaut2009mapping}
Pieter Blignaut. 2009.
\newblock Fixing the error: Mapping of eye-tracking data.
\newblock In \emph{Proceedings of the 2009 Symposium on Eye Tracking Research \& Applications (ETRA)}, 251--258.

\bibitem{blignaut2014mapping}
Pieter Blignaut. 2014.
\newblock Mapping the eye gaze: The impact of recalibration and head movements.
\newblock \emph{Journal of Eye Movement Research} 7, 3 (2014), 1--14.

\bibitem{lugaresi2019mediapipe}
Camillo Lugaresi, Jiuqiang Tang, Hadon Nash, Chris McClanahan, Esha Uboweja, Michael Hays, Fan Zhang, Chuo-Ling Chang, Ming~Guang Yong, Juhyun Lee, et al. 2019.
\newblock Mediapipe: A framework for building perception pipelines.
\newblock \emph{arXiv preprint arXiv:1906.08172}.

\bibitem{zou2005regularization}
Hui Zou and Trevor Hastie. 2005.
\newblock Regularization and variable selection via the elastic net.
\newblock \emph{Journal of the Royal Statistical Society: Series B (Statistical Methodology)} 67, 2 (2005), 301--320.

\bibitem{salvucci2000identifying}
Dario~D. Salvucci and Joseph~H. Goldberg. 2000.
\newblock Identifying fixations and saccades in eye-tracking protocols.
\newblock In \emph{Proceedings of the 2000 Symposium on Eye Tracking Research \& Applications (ETRA)}, 71--78.

\bibitem{hansen2010eye}
Dan~Witzner Hansen and Qiang Ji. 2010.
\newblock In the eye of the beholder: A survey of models for eyes and gaze.
\newblock \emph{IEEE Transactions on Pattern Analysis and Machine Intelligence} 32, 3 (2010), 478--500.

\bibitem{sesin2008investigation}
Andres Sesin, Armando Barreto, and Malek Adjouadi. 2008.
\newblock Investigation of artificial neural network architectures for pupil-corneal reflection eye tracking.
\newblock In \emph{Proceedings of the 2008 Symposium on Eye Tracking Research \& Applications (ETRA)}, 219--222.

\bibitem{sugano2014learning}
Yusuke Sugano, Yasuyuki Matsushita, and Yoichi Sato. 2014.
\newblock Learning-by-synthesis for appearance-based 3D gaze estimation.
\newblock In \emph{Proceedings of the IEEE Conference on Computer Vision and Pattern Recognition (CVPR)}, 1821--1828.

\end{thebibliography}

\end{document}
